\clearpage
\begin{trapbox}
\textbf{C094\quad 易错清单与关键边界}
\begin{itemize}
\item \textbf{「焦距相同」就误判成「共焦」。}
必须两焦点\emph{位置}相同；只知道两条曲线的焦半距都等于 $c$
但中心不同，不足以套公式。例如把共焦椭圆平移一个单位，
焦半距仍为 $2$，但焦点位置改变，已不再共焦。

\item \textbf{误把两处 $a$ 当同一个。}
椭圆是 $a_E>c$，双曲线是 $a_H<c$；不能写成一个公共 $a$。
两种半轴的定义不同：椭圆 $b_E^2=a_E^2-c^2$，
双曲线 $b_H^2=c^2-a_H^2$，正负号也不同。

\item \textbf{只写出 $(X,Y)$ 一个交点。}
两曲线方程都关于坐标轴对称，且 $X,Y>0$；
所以四点是 $(X,Y),(X,-Y),(-X,Y),(-X,-Y)$，
不能错误写成仅 $(X,Y),(-X,-Y)$ 两点。

\item \textbf{把 $(\pm X,\pm Y)$ 理解成同号取法。}
这里两处符号\emph{相互独立}。同号只得到第一、第三象限，
会漏掉第二、第四象限。

\item \textbf{误认为有些非退化共焦组合只有零个或两个交点。}
只要 $a_E>c>a_H>0$，就有 $X^2>0,Y^2>0$；
因而恰有四个\emph{互异的实交点}。
两点合并只能在 $b_E=0$ 或 $b_H=0$ 等退化极限出现，
已经不满足题设。

\item \textbf{双曲线虚半轴平方写错符号。}
应是 $b_H^2=c^2-a_H^2$，不是 $a_H^2-c^2$。
错误会令 $Y^2<0$，产生「无实交点」的假结论。

\item \textbf{开平方不检验正负与分母。}
$c>0$ 且 $b_E,b_H>0$，才可写 $Y=b_Eb_H/c$ 为第一象限正值。
应先求 $x^2,y^2$，再分别开正负根。

\item \textbf{把「切线正交」当成「焦点张角 $90^\circ$」。}
对例1，交点处切线斜率是 $-1/2$ 与 $2$，
但焦角满足 $\cos\theta=3/5$。两种角位于不同的直线对。

\item \textbf{忽视坐标系要求。}
公式 $(\pm X,\pm Y)$ 默认共同中心为原点、焦点轴为 $x$ 轴。
若原坐标系中中心为 $(h,k)$、焦点轴仍水平，
应写成 $(h\pm X,k\pm Y)$；若焦点轴旋转，还须做坐标旋转。
\end{itemize}
\end{trapbox}

\clearpage
\begin{trapbox}
\textbf{最小反例与边界检查（帮助学生判断公式能否使用）。}

\textbf{反例1：同焦半距但不是共焦。}
\[
E:\frac{x^2}{16}+\frac{y^2}{12}=1,\qquad
H:\frac{(x-1)^2}{1}-\frac{y^2}{3}=1.
\]
\begin{center}
% Identical axis scales; right panel translates H one unit.
\begin{tikzpicture}[x=.43cm,y=.43cm,font=\small]
\foreach \shift/\dx/\title in {0/0/同一对焦点,14/1/仅焦半距相同} {
\begin{scope}[shift={(\shift,0)}]
\node at (0,4.6) {\title};
\draw[gray] (-4.8,0)--(5,0);
\draw[blue!70!black,thick] (0,0) ellipse (4 and {sqrt(12)});
\foreach \s in {-1,1} {
\draw[orange!80!black,thick,dashed,domain=-3.5:3.5,samples=80,variable=\t]
 plot ({\dx+\s*sqrt(1+\t*\t/3)},\t);}
\foreach \f in {-2,2} {
\fill[blue!70!black] (\f,0) circle (3pt);
\draw[orange!80!black,thick] ({\f+\dx},0) circle (5pt);
\node[orange!80!black,above=6pt] at ({\f+\dx},0) {$\pgfmathparse{int(\f+\dx)}\pgfmathresult$};}
\end{scope}}
\node[blue!70!black] at (0,-4.6) {两图实点：$(\pm2,0)$};
\node[orange!80!black] at (14,-4.6) {空圈：双曲线焦点};
\end{tikzpicture}


{\small 左图共焦；右图双曲线平移一单位，焦半距未变，焦点位置已变。}
\end{center}
两图均取 $c=2$、相同坐标比例；只允许左图套用四点公式。

\textbf{反例2：改变共焦曲线的方向不是自由选项。}
同一对不同焦点 $(-c,0),(c,0)$ 唯一决定焦点轴为水平轴。
把双曲线写成开口上下的 $y^2/a_H^2-x^2/b_H^2=1$
会使焦点位于 $y$ 轴，除非改变坐标与焦点位置，
否则\emph{不是同一对焦点}。

\textbf{边界1（椭圆趋向退化）：}
$a_E\downarrow c$ 时，$b_E\downarrow0$、$Y\downarrow0$；
但 $a_E=c$ 对应退化曲线，不在定理范围内。
\textbf{边界2（双曲线趋向退化）：}
$a_H\uparrow c$ 时，$b_H\downarrow0$、$Y\downarrow0$；
但 $a_H=c$ 同样不满足非退化条件。
\textbf{边界3（两类曲线丢失正常定义）：}
$c=0$ 或 $a_H=0$ 均不属于本结论的条件。

\textbf{考场自检（代回不用重算）。}
\[
\frac{X^2}{a_E^2}=\frac{a_H^2}{c^2},\qquad
\frac{Y^2}{b_E^2}=\frac{b_H^2}{c^2},\quad
\text{两者相加}=1;
\]
\[
\frac{X^2}{a_H^2}=\frac{a_E^2}{c^2},\qquad
\frac{Y^2}{b_H^2}=\frac{b_E^2}{c^2},\quad
\text{前者减后者}=1.
\]
\textbf{这两行既核对代数，也检验几何参数的正负。}
\end{trapbox}
